\section{Introducción}

El proyecto del genoma humano, comenzado a finales de la década de 1980, era promovido como el proyecto final para descubrir la "huella digital de la vida". Se pensaba que el proyecto estaría terminado para el 2005 sin embargo, lograr descubrir la "huella digital de la vida" para ese año era bastante cuestionable. Para ese entonces ya se tenían secuenciados algunos genomas de bacterias y levaduras, y se sabía que, de la función biológica de los genes, entre el 30\% al 50\% permanecían aun sin caracterizar. A todo esto había que sumar que los genes y sus productos eran solo componentes individuales de un sistema biológico entero, el cómo se comporta cada uno de estos componentes no es suficiente para entender el sistema completo \parencite{Kanehisa_1997}.

Es por esto que surge la necesidad de crear una plataforma para efectuar el post-análisis de toda esta nueva información,  es decir realizar acercamientos experimentales e informáticos para sistemáticamente descubrir las interacciones y rutas génicas y metabólicas. Así surge la enciclopedia Kyoto de genes y genomas (KEGG), siendo una plataforma en línea para llevar a cabo este post-análisis \parencite{Kanehisa_1997}. KEGG partió teniendo tres objetivos:
\begin{enumerate}
	\item Computarizar el conocimiento actual de las rutas génicas y metabólicas de observaciones experimentales.
	\item Mantener un catálogo de genes para cada organismo que ha sido secuenciado, y cada componente del catálogo que sea mapeado a las rutas de KEGG.
	\item Desarrollar nuevas tecnologías informáticas que estén asociadas con rutas e interacciones. 
\end{enumerate}

La disponibilidad de genomas completos provenientes de distintos organismos para su análisis comparativo provee una nueva perspectiva sobre las relaciones homólogas entre los genes. Con toda esta nueva información sumada a todas las técnicas de análisis computacionales, \textcite{Koonin_2005} plantea que debería ser posible entonces reconstruir la historia evolutiva de cada gen (siempre y cuando se tenga secuenciado todo el genoma que lo contenga). Y a su vez ha tenido el efecto de reintroducir un significado mas amplio de familia génica para incluir a los genes parálogos dentro de una especie y genes ortólogos o parálogos entre distintas especies. 

La homología, en su definición más general, designa una relación de descendencia común entre cualquier entidad, sin alguna especificación adicional sobre el escenario evolutivo. De acuerdo con esto, las entidades estarían relacionadas por homología, si hablamos de genes, estos serían llamados homólogos \parencite{Koonin_2005}. 

El estudio y propuesta de grupos ortólogos ha sido una de las prácticas más frecuentes por los biólogos moleculares, incluso antes del desarrollo de la genómica comparada, análisis y estudio de las familias génicas \parencite{doctor}.

Las familias génicas son grupos de genes que descienden de un ancestro en común y retienen cierta similitud en su secuencia y a menudo tienen una función compartida \parencite{Dayhoff_1976}. El concepto de  familia génica aplica tanto a ambos genes dentro de un genoma (parálogos) y genes relacionados entre genomas (ortólogos y parálogos) \parencite{Dayhoff_1976}.

